% ============================================================================
%  common.tex —— 校本课插图公共绘图代码
%  由同目录下的各图 *.tex 用 % ============================================================================
%  common.tex —— 校本课插图公共绘图代码
%  由同目录下的各图 *.tex 用 % ============================================================================
%  common.tex —— 校本课插图公共绘图代码
%  由同目录下的各图 *.tex 用 % ============================================================================
%  common.tex —— 校本课插图公共绘图代码
%  由同目录下的各图 *.tex 用 \input{common.tex} 引入。
%  改配色 / 改基本图形尺寸，只改这一个文件即可，所有图同步生效。
%
%  编译：见同目录 README.md
% ============================================================================

% ---------- 颜色（只用于两件事：标出重点、区分正反面）----------
\definecolor{tuline}{RGB}{85,85,85}      % 深灰：普通线条；圆盘（表盘）的外圈与数字
\definecolor{tuhigh}{RGB}{214,126,0}     % 橙色：标出重点（两次演示的终点卡）
\definecolor{tuback}{RGB}{170,50,50}     % 深红：背面（翻过来的那一面）
\definecolor{tubfill}{RGB}{247,225,225}  % 浅红：背面填充
\definecolor{tugrid}{RGB}{150,150,150}   % 表格细线
\definecolor{tuhead}{RGB}{236,236,236}   % 表头底色

% ---------- 统一线宽 / 字号 ----------
% 画出来的图最终会被 \includegraphics[width=8cm] 等缩放，
% 下面的尺寸按“最终宽度 8cm 左右”设计，缩放比接近 1，字号损失很小。
\tikzset{
  tuline/.style={draw=tuline,   line width=0.8pt, line cap=round, line join=round},
  tuback/.style={draw=tuback,   line width=0.8pt, line cap=round, line join=round},
  tudot/.style={circle, fill=tuline,  inner sep=0pt, minimum size=1.7pt},
  tuarrb/.style={draw=tuline, line width=1.0pt, line cap=round, line join=round,
                -{Stealth[length=2.2mm,width=1.7mm]}},
}
\newcommand{\tufont}{\fontsize{13}{15.6}\selectfont}   % 图中主要文字
\newcommand{\tufonts}{\fontsize{10}{12}\selectfont}    % 图中小字（动作名、说明）
\newcommand{\tufontb}{\fontsize{17}{20}\selectfont}    % 表格里的算式
\newcommand{\tufonty}{\fontsize{16}{19.2}\selectfont}  % 圆盘（表盘）上的数字，要比正文字大

% ============================================================================
%  一、12 格圆盘 = 空表盘（图 1 / 2 / 3 共用同一段代码，保证逐像素一致）
%      长得像时钟，只有两样东西：外圈 + 一圈数字。
%      **不要**刻度线、不要高亮扇形、不要外弧、不要箭头。
%      需要演示的地方（8+7、11+1 之类）上课时由老师在图上手画。
%
%      位置约定（与课件、讲义一致）：0 在正上方（12 点），1,2,...,11 顺时针。
%      TikZ 的 0° 指向右、正角逆时针，所以第 k 个数字的方向角 = 90-30k 度：
%        k=0 → 90°（正上方）；k=3 → 0°（正右）；k=6 → -90°（正下方）；k=9 → 180°（正左）。
% ============================================================================
% 尺寸（改这两个数即可同时改三张圆盘；单位 cm）
\newcommand{\yuanpanR}{3.30}   % 外圈半径（图幅 = 2\yuanpanR + 一点边框）
\newcommand{\yuanpanM}{2.05}   % 数字所在半径（外圈以内，数字既不碰外圈也不挤在一起）

% 圆盘本体：外圈 + 一圈数字（空表盘，留给老师上课手画）
\newcommand{\yuanpanbase}{%
  \draw[tuline] (0,0) circle (\yuanpanR);
  \foreach \k in {0,...,11} {%
    \node at (90-\k*30:\yuanpanM) {\tufonty \(\k\)};
  }%
}

% ============================================================================
%  二、正三角形纸片（A 上、B 右下、C 左下），全套统一尺寸
%      图形自带一个 scope：把“中心”平移 (0,0) 的原点，再整体 scale=\triS
%      （内部坐标：circumradius = 1）
% ============================================================================
\newcommand{\triS}{1.58}          % 缩放：等边三角形边长 = 1.732*1.58 ≈ 2.74cm
% #1#2#3 = 上(A)、右下(B)、左下(C) 三个顶点上依次写的标签
\newcommand{\tricard}[3]{%
  \begin{scope}[scale=\triS]
    \draw[tuline] (0,1) -- (-0.86603,-0.5) -- (0.86603,-0.5) -- cycle;
    \node at (0,0.55)      {\tufont #1};%
    \node at (0.49,-0.29)  {\tufont #2};%
    \node at (-0.49,-0.29) {\tufont #3};%
  \end{scope}%
}

% ============================================================================
%  三、正方形纸片（1 左上、2 右上、3 右下、4 左下），全套统一尺寸
%      \squcard 的 scope 内：中心在原点，半边长 1，整体 scale=\sqS
% ============================================================================
\newcommand{\sqS}{1.24}           % 缩放：正方形边长 = 2*1.24 = 2.48cm
% #1#2#3#4 = 左上、右上、右下、左下 四个顶点上依次写的标签
\newcommand{\squcard}[4]{%
  \begin{scope}[scale=\sqS]
    \draw[tuline] (-1,1) -- (1,1) -- (1,-1) -- (-1,-1) -- cycle;
    \node at (-0.52,0.52)  {\tufont #1};%
    \node at (0.52,0.52)   {\tufont #2};%
    \node at (0.52,-0.52)  {\tufont #3};%
    \node at (-0.52,-0.52) {\tufont #4};%
  \end{scope}%
}
 引入。
%  改配色 / 改基本图形尺寸，只改这一个文件即可，所有图同步生效。
%
%  编译：见同目录 README.md
% ============================================================================

% ---------- 颜色（只用于两件事：标出重点、区分正反面）----------
\definecolor{tuline}{RGB}{85,85,85}      % 深灰：普通线条；圆盘（表盘）的外圈与数字
\definecolor{tuhigh}{RGB}{214,126,0}     % 橙色：标出重点（两次演示的终点卡）
\definecolor{tuback}{RGB}{170,50,50}     % 深红：背面（翻过来的那一面）
\definecolor{tubfill}{RGB}{247,225,225}  % 浅红：背面填充
\definecolor{tugrid}{RGB}{150,150,150}   % 表格细线
\definecolor{tuhead}{RGB}{236,236,236}   % 表头底色

% ---------- 统一线宽 / 字号 ----------
% 画出来的图最终会被 \includegraphics[width=8cm] 等缩放，
% 下面的尺寸按“最终宽度 8cm 左右”设计，缩放比接近 1，字号损失很小。
\tikzset{
  tuline/.style={draw=tuline,   line width=0.8pt, line cap=round, line join=round},
  tuback/.style={draw=tuback,   line width=0.8pt, line cap=round, line join=round},
  tudot/.style={circle, fill=tuline,  inner sep=0pt, minimum size=1.7pt},
  tuarrb/.style={draw=tuline, line width=1.0pt, line cap=round, line join=round,
                -{Stealth[length=2.2mm,width=1.7mm]}},
}
\newcommand{\tufont}{\fontsize{13}{15.6}\selectfont}   % 图中主要文字
\newcommand{\tufonts}{\fontsize{10}{12}\selectfont}    % 图中小字（动作名、说明）
\newcommand{\tufontb}{\fontsize{17}{20}\selectfont}    % 表格里的算式
\newcommand{\tufonty}{\fontsize{16}{19.2}\selectfont}  % 圆盘（表盘）上的数字，要比正文字大

% ============================================================================
%  一、12 格圆盘 = 空表盘（图 1 / 2 / 3 共用同一段代码，保证逐像素一致）
%      长得像时钟，只有两样东西：外圈 + 一圈数字。
%      **不要**刻度线、不要高亮扇形、不要外弧、不要箭头。
%      需要演示的地方（8+7、11+1 之类）上课时由老师在图上手画。
%
%      位置约定（与课件、讲义一致）：0 在正上方（12 点），1,2,...,11 顺时针。
%      TikZ 的 0° 指向右、正角逆时针，所以第 k 个数字的方向角 = 90-30k 度：
%        k=0 → 90°（正上方）；k=3 → 0°（正右）；k=6 → -90°（正下方）；k=9 → 180°（正左）。
% ============================================================================
% 尺寸（改这两个数即可同时改三张圆盘；单位 cm）
\newcommand{\yuanpanR}{3.30}   % 外圈半径（图幅 = 2\yuanpanR + 一点边框）
\newcommand{\yuanpanM}{2.05}   % 数字所在半径（外圈以内，数字既不碰外圈也不挤在一起）

% 圆盘本体：外圈 + 一圈数字（空表盘，留给老师上课手画）
\newcommand{\yuanpanbase}{%
  \draw[tuline] (0,0) circle (\yuanpanR);
  \foreach \k in {0,...,11} {%
    \node at (90-\k*30:\yuanpanM) {\tufonty \(\k\)};
  }%
}

% ============================================================================
%  二、正三角形纸片（A 上、B 右下、C 左下），全套统一尺寸
%      图形自带一个 scope：把“中心”平移 (0,0) 的原点，再整体 scale=\triS
%      （内部坐标：circumradius = 1）
% ============================================================================
\newcommand{\triS}{1.58}          % 缩放：等边三角形边长 = 1.732*1.58 ≈ 2.74cm
% #1#2#3 = 上(A)、右下(B)、左下(C) 三个顶点上依次写的标签
\newcommand{\tricard}[3]{%
  \begin{scope}[scale=\triS]
    \draw[tuline] (0,1) -- (-0.86603,-0.5) -- (0.86603,-0.5) -- cycle;
    \node at (0,0.55)      {\tufont #1};%
    \node at (0.49,-0.29)  {\tufont #2};%
    \node at (-0.49,-0.29) {\tufont #3};%
  \end{scope}%
}

% ============================================================================
%  三、正方形纸片（1 左上、2 右上、3 右下、4 左下），全套统一尺寸
%      \squcard 的 scope 内：中心在原点，半边长 1，整体 scale=\sqS
% ============================================================================
\newcommand{\sqS}{1.24}           % 缩放：正方形边长 = 2*1.24 = 2.48cm
% #1#2#3#4 = 左上、右上、右下、左下 四个顶点上依次写的标签
\newcommand{\squcard}[4]{%
  \begin{scope}[scale=\sqS]
    \draw[tuline] (-1,1) -- (1,1) -- (1,-1) -- (-1,-1) -- cycle;
    \node at (-0.52,0.52)  {\tufont #1};%
    \node at (0.52,0.52)   {\tufont #2};%
    \node at (0.52,-0.52)  {\tufont #3};%
    \node at (-0.52,-0.52) {\tufont #4};%
  \end{scope}%
}
 引入。
%  改配色 / 改基本图形尺寸，只改这一个文件即可，所有图同步生效。
%
%  编译：见同目录 README.md
% ============================================================================

% ---------- 颜色（只用于两件事：标出重点、区分正反面）----------
\definecolor{tuline}{RGB}{85,85,85}      % 深灰：普通线条；圆盘（表盘）的外圈与数字
\definecolor{tuhigh}{RGB}{214,126,0}     % 橙色：标出重点（两次演示的终点卡）
\definecolor{tuback}{RGB}{170,50,50}     % 深红：背面（翻过来的那一面）
\definecolor{tubfill}{RGB}{247,225,225}  % 浅红：背面填充
\definecolor{tugrid}{RGB}{150,150,150}   % 表格细线
\definecolor{tuhead}{RGB}{236,236,236}   % 表头底色

% ---------- 统一线宽 / 字号 ----------
% 画出来的图最终会被 \includegraphics[width=8cm] 等缩放，
% 下面的尺寸按“最终宽度 8cm 左右”设计，缩放比接近 1，字号损失很小。
\tikzset{
  tuline/.style={draw=tuline,   line width=0.8pt, line cap=round, line join=round},
  tuback/.style={draw=tuback,   line width=0.8pt, line cap=round, line join=round},
  tudot/.style={circle, fill=tuline,  inner sep=0pt, minimum size=1.7pt},
  tuarrb/.style={draw=tuline, line width=1.0pt, line cap=round, line join=round,
                -{Stealth[length=2.2mm,width=1.7mm]}},
}
\newcommand{\tufont}{\fontsize{13}{15.6}\selectfont}   % 图中主要文字
\newcommand{\tufonts}{\fontsize{10}{12}\selectfont}    % 图中小字（动作名、说明）
\newcommand{\tufontb}{\fontsize{17}{20}\selectfont}    % 表格里的算式
\newcommand{\tufonty}{\fontsize{16}{19.2}\selectfont}  % 圆盘（表盘）上的数字，要比正文字大

% ============================================================================
%  一、12 格圆盘 = 空表盘（图 1 / 2 / 3 共用同一段代码，保证逐像素一致）
%      长得像时钟，只有两样东西：外圈 + 一圈数字。
%      **不要**刻度线、不要高亮扇形、不要外弧、不要箭头。
%      需要演示的地方（8+7、11+1 之类）上课时由老师在图上手画。
%
%      位置约定（与课件、讲义一致）：0 在正上方（12 点），1,2,...,11 顺时针。
%      TikZ 的 0° 指向右、正角逆时针，所以第 k 个数字的方向角 = 90-30k 度：
%        k=0 → 90°（正上方）；k=3 → 0°（正右）；k=6 → -90°（正下方）；k=9 → 180°（正左）。
% ============================================================================
% 尺寸（改这两个数即可同时改三张圆盘；单位 cm）
\newcommand{\yuanpanR}{3.30}   % 外圈半径（图幅 = 2\yuanpanR + 一点边框）
\newcommand{\yuanpanM}{2.05}   % 数字所在半径（外圈以内，数字既不碰外圈也不挤在一起）

% 圆盘本体：外圈 + 一圈数字（空表盘，留给老师上课手画）
\newcommand{\yuanpanbase}{%
  \draw[tuline] (0,0) circle (\yuanpanR);
  \foreach \k in {0,...,11} {%
    \node at (90-\k*30:\yuanpanM) {\tufonty \(\k\)};
  }%
}

% ============================================================================
%  二、正三角形纸片（A 上、B 右下、C 左下），全套统一尺寸
%      图形自带一个 scope：把“中心”平移 (0,0) 的原点，再整体 scale=\triS
%      （内部坐标：circumradius = 1）
% ============================================================================
\newcommand{\triS}{1.58}          % 缩放：等边三角形边长 = 1.732*1.58 ≈ 2.74cm
% #1#2#3 = 上(A)、右下(B)、左下(C) 三个顶点上依次写的标签
\newcommand{\tricard}[3]{%
  \begin{scope}[scale=\triS]
    \draw[tuline] (0,1) -- (-0.86603,-0.5) -- (0.86603,-0.5) -- cycle;
    \node at (0,0.55)      {\tufont #1};%
    \node at (0.49,-0.29)  {\tufont #2};%
    \node at (-0.49,-0.29) {\tufont #3};%
  \end{scope}%
}

% ============================================================================
%  三、正方形纸片（1 左上、2 右上、3 右下、4 左下），全套统一尺寸
%      \squcard 的 scope 内：中心在原点，半边长 1，整体 scale=\sqS
% ============================================================================
\newcommand{\sqS}{1.24}           % 缩放：正方形边长 = 2*1.24 = 2.48cm
% #1#2#3#4 = 左上、右上、右下、左下 四个顶点上依次写的标签
\newcommand{\squcard}[4]{%
  \begin{scope}[scale=\sqS]
    \draw[tuline] (-1,1) -- (1,1) -- (1,-1) -- (-1,-1) -- cycle;
    \node at (-0.52,0.52)  {\tufont #1};%
    \node at (0.52,0.52)   {\tufont #2};%
    \node at (0.52,-0.52)  {\tufont #3};%
    \node at (-0.52,-0.52) {\tufont #4};%
  \end{scope}%
}
 引入。
%  改配色 / 改基本图形尺寸，只改这一个文件即可，所有图同步生效。
%
%  编译：见同目录 README.md
% ============================================================================

% ---------- 颜色（只用于两件事：标出重点、区分正反面）----------
\definecolor{tuline}{RGB}{85,85,85}      % 深灰：普通线条；圆盘（表盘）的外圈与数字
\definecolor{tuhigh}{RGB}{214,126,0}     % 橙色：标出重点（两次演示的终点卡）
\definecolor{tuback}{RGB}{170,50,50}     % 深红：背面（翻过来的那一面）
\definecolor{tubfill}{RGB}{247,225,225}  % 浅红：背面填充
\definecolor{tugrid}{RGB}{150,150,150}   % 表格细线
\definecolor{tuhead}{RGB}{236,236,236}   % 表头底色

% ---------- 统一线宽 / 字号 ----------
% 画出来的图最终会被 \includegraphics[width=8cm] 等缩放，
% 下面的尺寸按“最终宽度 8cm 左右”设计，缩放比接近 1，字号损失很小。
\tikzset{
  tuline/.style={draw=tuline,   line width=0.8pt, line cap=round, line join=round},
  tuback/.style={draw=tuback,   line width=0.8pt, line cap=round, line join=round},
  tudot/.style={circle, fill=tuline,  inner sep=0pt, minimum size=1.7pt},
  tuarrb/.style={draw=tuline, line width=1.0pt, line cap=round, line join=round,
                -{Stealth[length=2.2mm,width=1.7mm]}},
}
\newcommand{\tufont}{\fontsize{13}{15.6}\selectfont}   % 图中主要文字
\newcommand{\tufonts}{\fontsize{10}{12}\selectfont}    % 图中小字（动作名、说明）
\newcommand{\tufontb}{\fontsize{17}{20}\selectfont}    % 表格里的算式
\newcommand{\tufonty}{\fontsize{16}{19.2}\selectfont}  % 圆盘（表盘）上的数字，要比正文字大

% ============================================================================
%  一、12 格圆盘 = 空表盘（图 1 / 2 / 3 共用同一段代码，保证逐像素一致）
%      长得像时钟，只有两样东西：外圈 + 一圈数字。
%      **不要**刻度线、不要高亮扇形、不要外弧、不要箭头。
%      需要演示的地方（8+7、11+1 之类）上课时由老师在图上手画。
%
%      位置约定（与课件、讲义一致）：0 在正上方（12 点），1,2,...,11 顺时针。
%      TikZ 的 0° 指向右、正角逆时针，所以第 k 个数字的方向角 = 90-30k 度：
%        k=0 → 90°（正上方）；k=3 → 0°（正右）；k=6 → -90°（正下方）；k=9 → 180°（正左）。
% ============================================================================
% 尺寸（改这两个数即可同时改三张圆盘；单位 cm）
\newcommand{\yuanpanR}{3.30}   % 外圈半径（图幅 = 2\yuanpanR + 一点边框）
\newcommand{\yuanpanM}{2.05}   % 数字所在半径（外圈以内，数字既不碰外圈也不挤在一起）

% 圆盘本体：外圈 + 一圈数字（空表盘，留给老师上课手画）
\newcommand{\yuanpanbase}{%
  \draw[tuline] (0,0) circle (\yuanpanR);
  \foreach \k in {0,...,11} {%
    \node at (90-\k*30:\yuanpanM) {\tufonty \(\k\)};
  }%
}

% ============================================================================
%  二、正三角形纸片（A 上、B 右下、C 左下），全套统一尺寸
%      图形自带一个 scope：把“中心”平移 (0,0) 的原点，再整体 scale=\triS
%      （内部坐标：circumradius = 1）
% ============================================================================
\newcommand{\triS}{1.58}          % 缩放：等边三角形边长 = 1.732*1.58 ≈ 2.74cm
% #1#2#3 = 上(A)、右下(B)、左下(C) 三个顶点上依次写的标签
\newcommand{\tricard}[3]{%
  \begin{scope}[scale=\triS]
    \draw[tuline] (0,1) -- (-0.86603,-0.5) -- (0.86603,-0.5) -- cycle;
    \node at (0,0.55)      {\tufont #1};%
    \node at (0.49,-0.29)  {\tufont #2};%
    \node at (-0.49,-0.29) {\tufont #3};%
  \end{scope}%
}

% ============================================================================
%  三、正方形纸片（1 左上、2 右上、3 右下、4 左下），全套统一尺寸
%      \squcard 的 scope 内：中心在原点，半边长 1，整体 scale=\sqS
% ============================================================================
\newcommand{\sqS}{1.24}           % 缩放：正方形边长 = 2*1.24 = 2.48cm
% #1#2#3#4 = 左上、右上、右下、左下 四个顶点上依次写的标签
\newcommand{\squcard}[4]{%
  \begin{scope}[scale=\sqS]
    \draw[tuline] (-1,1) -- (1,1) -- (1,-1) -- (-1,-1) -- cycle;
    \node at (-0.52,0.52)  {\tufont #1};%
    \node at (0.52,0.52)   {\tufont #2};%
    \node at (0.52,-0.52)  {\tufont #3};%
    \node at (-0.52,-0.52) {\tufont #4};%
  \end{scope}%
}
